% Chapter 9: Miniatures of India and Iran: The System of Free Perceptual Perspective
% Source pages: 190–214

\chapter{Miniatures of India and Iran: The System of Free Perceptual Perspective}

\markboth{Chapter IX. Miniatures of India and Iran}{Chapter IX}

\srcpage{190}

The study of the spatial constructions characteristic of medieval Indian and Iranian
miniatures is limited here to an attempt to discover in them the general geometric
regularities of medieval painting that were investigated in earlier chapters with
reference to Byzantine and Old Russian works of
art.\footnote{The author does not aim to clarify datings or attributions either.
  The data used below are therefore taken from the following sources: Armenag Bey
  Sakisian. \emph{La miniature persane du XII au XVII si\`ecle}. Paris, 1929
  (ills.\,87, 89, 92, 94, 95, 96); J.\,Barrett, B.\,Gray. \emph{La peinture
  indienne}. Gen\`eve, 1963 (ills.\,86, 88); B.\,W.\,Robinson. \emph{Persian
  Painting in the Metropolitan Museum of Art}. New York, 1953 (ill.\,90);
  B.\,Gray. \emph{Rajput Painting}. London, 1948 (ill.\,93); B.\,Gray.
  \emph{Persische Malerei}. Gen\`eve, 1961 (ill.\,97); M.\,M.\,Ashrafi.
  \emph{Persian-Tajik Poetry in Miniatures of the 14th--17th Centuries}.
  Dushanbe, 1974 (ill.\,91). The originals' locations are also given in these
  sources.}

The aim pursued here makes a historical approach to the analysis of miniatures
unnecessary. The complex and many-layered influence of different artistic schools
on one another, the various stages in the historical development of style, and
analogous problems all lie outside the scope of the analysis conducted in this
chapter. What matters for us here is to identify the rational foundations of the
distinctive methods of spatial construction employed with such success by the
miniaturists of India and Iran. For this reason, the illustrations drawn upon below
will be selected according to geometric rather than historical logic.

Before approaching a consideration of the miniatures, the first question to answer
is which space---objective or perceptual---they depict. Even a cursory analysis
reveals that both may be found, and often a synthesis of the two approaches,
analogous to what is observed in medieval European art.

As an example of an Indian miniature that conveys objective space, consider the one
from the pre-Mughal period (ill.\,86) (Malwa, c.\,1530). Everything here speaks of
a clearly expressed desire to convey the geometry of objective space: the
architecture rendered in the manner of a technical drawing, the absence of any
depiction of objectively horizontal surfaces, and the placement of figures' feet on
supporting lines (all of which are, as we recall, signs of orthographic projection).

\figblock{86}{Indian miniature. Malwa. c.\,1530.}

Noteworthy is the geometric closeness of this miniature to ancient Egyptian
painting, extending even to the depiction of the human figure---the feet and heads
shown in profile, the shoulders facing front. Since direct borrowing is out of the
question here, one can only conclude that the ancient Egyptian method of depicting
the human figure is natural and in some sense optimal, since Indian miniaturists,
confronted more than two thousand years later with the same task---to convey
objective space---arrived at the same solution. It should not be supposed that the
conveyance of objective-space geometry is a peculiarity of pre-Mughal miniatures.
The miniature of the Baghdad school (14th century) (ill.\,87), although it renders
human figures more freely (somewhat resembling those in ill.\,10), must also be
classified as a technical drawing rather than a pictorial rendering (Baghdad does
not belong to India or Iran, but its inclusion is justified by the close kinship of
the artistic schools of those regions).

Perceptual space may equally be found. A miniature in the provincial Mughal style
(c.\,1615--1620) depicts a courtly scene (ill.\,88). Here axonometry is used to
convey the geometry of perceptual space---which is entirely natural for a close and
shallow space.

\srcpage{192}

In the miniatures under discussion, the dominant depiction is of perceptual space,
with drawing methods serving only as auxiliaries. As in European medieval art,
axonometry forms the perspective foundation of painting here as well. This is
connected, on the one hand, with the fact that the artist frequently depicts
comparatively close and shallow layers of space, and on the other with the working
practice of the artist: he created his miniatures in enclosed spaces, from memory,
without recourse to painting from life. Even when depicting events taking place,
say, in a garden, he rendered every individual object as it would appear indoors or
at least from a sufficiently close distance. The axonometric basis of the painting
is also evident in the equal size of the figures depicted in the miniature. It
manifests itself, moreover, in the problem of limiting space, which was discussed
at length in Chapter~IV in connection with Byzantine and Old Russian painting.

\figblock{87}{Miniature. Baghdad school, 14th century.}

\figblock{88}{Courtly scene. Provincial Mughal style. c.\,1615--1620.}

The Timurid-style miniature (1410) (ill.\,89) is typical in all these respects.
The carpet and throne are rendered in axonometry, all the figures are of equal
size, and the space is bounded by a rocky hill depicted behind the throne, which
makes it unnecessary to show the far distance. Some figures as it were ``peek''
from behind the hill (a device frequently employed in miniatures), while the
treetops (of trees growing just behind the hill) remove the need to depict more
distant areas of space.

\srcpage{193}

Unlike European medieval painting, the depiction of the floor in interiors---or of
the ground surface when the action takes place in a courtyard, garden, and the
like---plays an enormous role in the miniatures under discussion. This was
stimulated by the fact that human life in the East is, so to speak, ``tied'' to the
floor. The floor is frequently adorned with beautiful carpets competing with the
decorative furnishings of the walls. All of this draws the artist's attention to the
floor and to the task of rendering it expressively in the miniature. That this task
is far from elementary was demonstrated by Chapters~III and~VI. Fig.\,23B showed
the difficulties of reconciling the local axonometries of the central section of an
interior's floor with those of the floor sections adjacent to the side walls.
Figs.\,46 and~47 showed how this reconciliation (in connection with ceiling
depiction) was handled in antiquity and at the close of the Western European Middle
Ages. It should be noted that European methods are fundamentally incapable of
solving the problems of the Eastern miniature, since they in practice amount to
distorting the depiction in the main, central part of the composition (a break in
the ceiling, the rendering of structural forms that do not really exist). It was
precisely these circumstances that pushed Eastern artists toward the direct use of
the scheme shown in Fig.\,23B. They accepted the visually absurd ``joining'' of the
central and lateral axonometries in order to render these three locally axonometric
regions as accurately as possible.

\figblock{89}{The Court of Khosrow. Shiraz. Timurid style. 1410.}

\srcpage{196}

One of the miniatures (Herat, c.\,1430) (ill.\,90) gives a typical depiction of an
interior floor and its walls. In the central section of the miniature, a series of
vertical and horizontal lines conveying the structure of the floor adjacent to the
far wall transitions ``discontinuously'' into two series of vertical and inclined
lines parallel to the bases of the side walls of the interior. An analogous
structure is used for the carpets (or divans) adjacent to the far and side walls,
and the objects placed on these are rendered in axonometry, confirming the
perspective approach to spatial representation in this miniature. Frequently, in
order to avoid the ``jump'' described above, artists extend the
vertical-horizontal structure of the central floor section all the way to the side
walls, accepting that the sections immediately adjacent to the side walls will
suffer distortion.

The vertical-horizontal floor structure shown here is remarkable in two respects.
First, it accurately conveys the visual perception of the middle section of such a
floor---though only for a narrow zone of clear central vision. The artist (going
against visual perception) extends this mode of depiction to left and right, wishing
to preserve in the miniature the parallelism of the lines that exists in objective
space. Second, such a depiction simultaneously constitutes a plan---a
technical-drawing rendering of objective-space geometry. Analysis of miniatures
shows that the artist always senses this dual character of such a depiction. The
floor in ill.\,90 cannot be entirely classified as a plan, since a plan would not
permit depicting the side walls; but the central section can be interpreted both as
a plan, and the miniaturist frequently takes advantage of this. The existence among
the miniatures of completely technical-drawing works (ills.\,86, 87) testifies that
this approach was perfectly comprehensible to both artist and viewer.

\figblock{90}{Bahram Gur in the Black Tower. Herat. c.\,1430.}

\srcpage{197}

The technical-drawing plan approach is especially clearly expressed in those cases
where, for example, a pool and the stone-lined channels feeding it are depicted in
the foreground of the miniature. In the Tabriz-style miniature (1524--1525)
(ill.\,91), the pavilion where the ruler is seated is rendered in clear axonometry,
while the stone-lined ground surface before him is shown in vertical-horizontal
structure, which can be understood both as conveying visual perception and as a
plan. The latter interpretation is confirmed by the depiction of a square pool with
its feeding channels in top view, on whose stone border a vase of flowers is shown.
Such a depiction of the vase (in a conditional right-angle rotation) is possible
only in a technical drawing---as was discussed in detail in Chapter~I. Thus, in
that miniature the symbiosis of pictorial rendering and technical drawing is
unambiguously apparent. Conditional right-angle rotations (an exact side view) for
objects standing on the ground (shown in plan) are a drawing device frequently used
by miniaturists.

\figblock{91}{Audience with the ruler. Tabriz style. 1524--1525. (tracing)}

\srcpage{198}

\figblock{92}{Small table with vessels (tracing). Early 14th century. Fragment.}

Carpets lying on the floor (or ground) are depicted by a similar device. As a rule
they are rendered in plan, which brings their depiction close to that of a technical
drawing, while simultaneously allowing the artist to preserve their multicoloured
patterns without perspective distortion. Several features also indicate that the
miniaturist followed visual perception: for small-format carpets this mode of seeing
is accurate. Evidence for this is that objects placed on the carpet are frequently
not shown in strict side view as they would be in a drawing, and that carpets
situated at the left or right margin of the miniature are shown in
foreshortening---in accordance with the scheme shown in Fig.\,23B.

The technical-drawing approach is not, however, dominant, and aside from the
relatively few fully technical-drawing miniatures (of the kind shown in ills.\,86,
87), it is confined to the depiction of floors and ground surfaces---and even there
it is often understood dually, as both a plan rendering and a rendering of natural
visual perception. This is evident in particular in occurrences of objects placed on
a single carpet being shown simultaneously in axonometry and in strict technical
side view.

As in Western medieval art, axonometry forms only the perspective basis of
depiction, so let us now examine the transformations characteristic of it. In the
foreground of a Mongol miniature of the early 14th century, a small table with
vessels is shown (ill.\,92). It is rendered taking account of the form-constancy
mechanism, according to scheme~40B. Unlike European artists, the Indian miniaturist
attaches no significance to the fact that the vessels standing on the table are
projected onto the surface of the ground. This is characteristic of the miniatures
under discussion, as confirmed in particular by the completely analogous mode of
rendering vessels standing on the carpet in the miniature considered above (ill.\,89).

\srcpage{199}

\figblock{93}{Genre scene. Pahari, Basoli style. 17th century.}

A different type of reverse-perspective construction can be seen in the Indian
miniature of the Basoli style depicting a genre scene (ill.\,93). Here the divan
is shown not according to scheme~40B but according to the rarely encountered
scheme~40C. The miniaturist preferred this for perfectly understandable reasons:
since the divan's legs are small, elongating them as scheme~B requires would make
the difference in height between the front and rear legs too unnatural, whereas the
objectively tall legs of the table in ill.\,92 allowed them in that case to remain
within visually reasonable bounds. Such a depiction of the surfaces of divans is
not exceptional; when they are rendered in perspective, their far edges are always
``raised'' in miniatures, as the form-constancy mechanism demands. As for the divan
surface itself, in the miniature under discussion it is shown in gentle reverse
perspective, corresponding to natural visual perception. The canopy corners,
likewise in accordance with the form-constancy mechanism, are rendered as almost
exact right angles---analogously to the way in which the defining silhouette angles
of Gospels were rendered in Byzantine art (ill.\,43). The floor is rendered in the
vertical linear structure already discussed above. Everything said here speaks of
the artist's effort to convey his visual perception of the nearby areas of space as
accurately as possible. The inevitable distortions in this case are shifted onto the
``detachment'' of the divan's rear legs from the floor. The sole violation of strict
adherence to visual perception---where not required by the rule of inevitable
distortions---concerns the depiction of the struts supporting the canopy: the lower
end of the front strut is shifted inward. Here recourse to the method of
geometrically contradictory depictions is entirely comprehensible, with absolutely
the same meaning as in the Byzantine miniature shown in ill.\,77.

\srcpage{200}

Sometimes in miniatures one encounters an unnaturally strong reverse perspective,
connected with the ``effects of reverse perspective'' discussed in detail in
Chapter~VI. The author of a miniature from the Herat school narrating the
adventures of Bahram Gur (1411) shows the left and right sides of an architectural
structure simultaneously (ill.\,94), producing the effect of strong reverse
perspective well known in Byzantine art (ill.\,54, right).

This miniature is interesting in another respect as well: technical drawing and
pictorial rendering are organically interwoven in it. Indeed, the floor rendered in
vertical-horizontal structure can be perceived both as its correct perspective
depiction (supported by the forms of objects placed on it) and as a plan (as
indicated by the depiction of a round pool in the very foreground). The surface of
the divan is turned ``toward the viewer'' (analogously to ill.\,93). At the same
time it is turned so strongly that one involuntarily perceives it as shown from
above, although this is contradicted by the depictions of human figures, which are
difficult to interpret as rendered in a conventional right-angle rotation. As a
result, horizontal surfaces (floor, divan) and vertical ones (far wall, side of
divan) as it were flatten out onto the miniature's surface. This approach is
unimpeachable from the standpoint of technical-drawing methods that permit
conventional rotations of projection planes, and can only be interpreted as a
depiction of perceptual space with certain caveats regarding perspective
exaggerations. Such duality and indeterminacy of geometry is usually quite justified
artistically: horizontal surfaces (carpets, for example) and vertical ones (walls)
alike are covered with expressive patterns and ornaments, and their undistorted
rendering is aesthetically well founded.

\srcpage{201}

The compatibility of vertical and horizontal planes of the subject with the picture
plane makes the ``decoding'' of the geometry of the depicted space not always
obvious at first glance. Let us turn first to a Herat-style miniature (1525)
(ill.\,95), which is analogous in many ways to the preceding one. The floor is
rendered similarly: of two objects standing on it, one is shown in axonometry and
the other in a conventional technical rotation; the carpet is flattened out, as the
divan surface was in the preceding miniature, leaving the viewer uncertain whether
the carpet is shown in natural visual perception or in plan. The balcony at upper
right is characterised by natural gentle reverse perspective. When one begins to
compare the carpet's depiction with the vertical wall carrying the balcony, it
sometimes seems as if that wall and the carpet are depictions of one and the same
plane---which is, of course, geometrically absurd. The only rational explanation
for what is depicted is the assumption that the carpet is located inside a room that
opens onto the wall with the balcony, and that this interior is rendered in a strict
orthographic projection that does not permit depiction of the side walls of the
interior.

An even stranger impression is produced by the combination of technical-drawing and
perspective devices in a Tabriz-style miniature depicting a scene of preaching
(first half of the 16th century) (ill.\,96). Here the floor and a series of carpets
with figures in conventional technical rotation are rendered in the scheme of a
``plan''. The stepped platform depicted in axonometry at the right creates the
impression that the figure shown at upper left is seated ``on the wall''---which is
incorrect, since his carpet extends only to the base of the far ornamental wall.
However, the combination of drawing and painting has made necessary in this
miniature the use of a geometrically contradictory depiction: the stepped platform
at the right is as it were leaning against the far wall at its top, while at its
base it is separated from the wall by approximately the width of the carpet shown
at upper left---geometrically analogous to the column-base shifts toward the viewer
described in Chapter~VIII (ill.\,76). This is indicated by the position of the
platform's legs, which is itself contradictory: the middle side leg touches the
floor not at the point that corresponds to the overall axonometric character of the
depiction.

\figblock{95}{Mahmud Muzahhib. Bahram Gur in the Yellow Tower. Herat style. 1525.}

\figblock{96}{Sheikh Zade. Scene of a Sermon. Tabriz style. First half of the 16th
  century.}

\srcpage{202}

\figblock{94}{Bahram Gur and one of the seven beauties. Herat style. 1411.}

In the Herat-style miniature (1494) (ill.\,97), depicting a scene of children
being taught, there is an entirely different type of geometrically contradictory
depiction: the right wall of the architectural structure shown on the second plane
conveys in its upper portion the exterior, and in its lower portion the interior of
the wall. In other respects the geometry of this miniature contains nothing new:
the ground and carpets are rendered as usual, with a tendency toward a plan
approach. The pool with its feeding channel is shown exactly in plan (and the
fountain in the centre of the pool in a strict conventional right-angle rotation).
Everything else is axonometry with a tendency toward gentle reverse perspective
(the fence beside the tree).

Something analogous in the matter of geometric contradiction is visible in another
miniature (ill.\,95) as well: the upper part of the wall with the dome conveys the
exterior, but this wall at the bottom becomes part of the interior. The wall
carrying the balcony is at the top positioned frontally and belongs to the exterior
of the building, while at the bottom it is shown in side foreshortening and relates
to the interior.

These devices of geometric contradiction make it possible to show the most
significant features of the interior and exterior of a building simultaneously, and
in addition allow the creation of expressive ``decorative frames'' around the main
figures. As for the transitions from the ``exterior'' to the ``interior'' view,
these are effected with remarkable tact and are entirely imperceptible, especially
since the artist sometimes considers it appropriate to provide ``camouflage'' for
the places that might betray the geometric contradiction. Precisely this function is
performed by the foliage of the tree depicted in the miniature (ill.\,97): it
removes the potential puzzling question about the connection between the interior
wall of the building and the exterior courtyard space in the lower right portion of
the depicted building. Such simultaneous depiction of the ``exterior'' and
``interior'' view is not foreign to Old Russian art either, where it is, however,
achieved on the basis of technical-drawing methods (cross-sections, see ill.\,67).

\srcpage{204}

\figblock{97}{Behzad. Layla and Majnun at School. Herat style. 1494.}

Although the miniatures of Iran and India generally convey deeper space than
Byzantine and Old Russian painting, this does not mean that the space they convey
is deep in the same sense as the space of modern-day art. Rather, what occupies the
artist in the miniatures under discussion is the problem of narrative
space---a space within whose framework he is able to tell the depicted story with
the desired degree of detail. The sensation of space is created, as is well known,
not by geometry alone but also by colour, composition, and other artistic means,
and if one returns with this in mind to the evaluation of the spatial character of
the miniatures, it is easy to note the traditional attachment of the artist to the
principle of planarity. Contributing to this are the axonometric basis of the
drawing, which prevents conveying the diminution of objects' apparent sizes with
increasing depth; the generous use of technical-drawing devices with the rendering
of horizontal surfaces in plan; and other means not discussed in this book, such as
the rhythmic structure of colour patches that holds the image in a single plane, the
carpet-like ornamentality of pictorial details and framings, the silhouette quality
of figures, and so on.

\srcpage{205}

The brief survey of the characteristic features of medieval Indian and Iranian
miniatures conducted above has shown that they employ the same elements of spatial
construction---local axonometry and its transformations toward reverse perspective,
technical-drawing methods, geometrically contradictory depictions---as Old Russian
and Byzantine painting. These devices are not merely similar; they are employed for
the same purposes. Consequently, the spatial constructions of the miniatures under
discussion are just as rationally founded as those in Old Russian painting. Moreover,
if medieval European, Chinese, Korean, and Japanese art were to be examined from
this standpoint, the same conclusion would remain valid.

Even in those cases (in later periods) where miniatures depict distant areas of
space with reduced depictions of objects in the far distance, the local-axonometric
approach is preserved---which accords naturally with the regularities of visual
perception, which becomes axonometric twice over: in close and in distant planes.
Connected with this is the fact that in the medieval art of the Far East, axonometry
is used not only for the rendering of distant areas of space but also of shallow
interiors.

All of this allows one to assert that the spatial constructions of medieval art in
different regions share a common geometric basis and common methods. Yet from
identical bricks one can build different structures---and paintings executed using
identical geometric methods may turn out to be entirely unlike one another.

\srcpage{206}

\begin{center}$*\quad*\quad*$\end{center}

\noindent Let us consider the organisation of pictorial space as a whole (its
perspective aspect only) in Old Russian and Byzantine art on the one hand, and in
Indian and Iranian miniature art on the other. This comparison will allow us to
identify certain characteristic features of the system of free perceptual
perspective. The geometric organisation of space borders directly on composition,
and certain compositional questions will therefore also be touched upon below. But
a whole range of geometric problems---in particular the problems of geometric
rhythm, symmetry, asymmetry, and analogous questions---will in practice be left
aside. What follows, then, concerns only those matters that can, with certain
reservations, be assigned to tasks connected with constructing a perspective system,
even if a free one, yet one subject to known requirements.

If the geometric structure of any rigid perspective system---even that of linear
perspective---exists independently of the content of a painting and can in principle
be identified in depictions of entirely neutral objects such as regular geometric
solids, then in the system of free perceptual perspective the artist's freedom
manifests itself not only in the depiction of individual objects (discussed in
detail above) but also in the organisation of the spatial structure of the painting
as a whole. The artist is entitled to vary this structure according to the artistic
tasks before him.

When one turns to depictions with a developed narrative, where it is appropriate to
speak of perspective constructions, the space in all the paintings under
consideration is divided into two planes. However, the role of these planes in
Byzantine and Old Russian painting and in Indian and Iranian miniatures is
diametrically opposed.

In Eastern miniatures the depicted figures live their own lives, without ``feeling''
that they are being observed. This life takes place frequently in a comparatively
large space in which the figures are capable of varied interaction. The viewer
watches from a distance, as it were; the pictorial narrative ignores him. Since the
space of action depicted in the miniature is comparatively large, the artist faces
the problem of representing a sufficiently extensive and deep interior, a section of
a courtyard, a garden, and so on. In this connection, working in the system of local
axonometries, the artist is constantly compelled to resort to scheme~23B, using it
fully or in part. This scheme is admirably suited to the depiction of events in
which the main figure is placed on the second plane, with the space before him given
over to the remaining participants. It thus becomes possible to ``direct'' the
viewer's attention to the main figure by placing the others to the left and right of
the compositional centre (ills.\,89, 90, 95, 96). Frequently the main figures are
as it were incorporated into a decorative frame of ornamental wall elements and
patterned carpets (ills.\,94, 88), and between them and the lower edge of the image
there is always a sufficiently extended layer of space.

\srcpage{208}

The geometric structure of space in Byzantine and Old Russian painting is entirely
different. This is because the main figures are now called to active interaction
with the viewer; the latter must in some sense be a participant in the event. These
main figures must be brought as close as possible to the beholder; nothing must
stand between them and the viewer to obstruct direct contact; they are painted,
so to speak, ``standing on the frame'', and moreover by an artist who has approached
the depicted figures very closely. But then the artist can only paint a group of
such figures one by one, each time as it were standing before the one being depicted.
The absence of space in front of them and the small distance from artist to depicted
figure makes it virtually impossible to paint such a group from a single
``organising'' viewpoint.

Under these conditions the geometry of pictorial space is built with the aid of the
second plane, which now plays a subordinate role. In Indian and Iranian miniatures
the lines of the side walls of the interior, the lateral carpets, the groups of
standing or seated people---all directed toward the semantic and geometric
centre---generally conform to the geometry of scheme~23B and as it were guide the
viewer's attention from the periphery of the first plane to the centre of the
second. In works of Old Russian and Byzantine art, the analogous emphasis of the
semantic and geometric centres can only run in the opposite direction---from the
periphery of the second plane to the centre of the first. Thus arises a geometric
structure bearing the character of reverse perspective, although it has (unlike
scheme~23B) not the slightest connection with the instantaneous visual perception
of space. As examples of such a spatial geometry, reference may be made to a number
of works (ills.\,27, 52, 68, 69, 82), in which the orientation of the architectural
structures of the second plane as it were focuses the viewer's attention on the
centre of the first plane.

Instead of highlighting the main figures with a decorative frame---here, since the
artist has free disposal only of the objects of the second plane---attention is drawn
to the main figures by depicting tall architectural structures immediately behind
them on the second plane (ills.\,28, 53), or conversely by a ``break'' in the tall
formations of the second plane (ill.\,35).

The subordinate character of the second plane in Byzantine and Old Russian painting
is evident from the fact that it is possible to compose the picture without it at
all. The iconography of the ``Crucifixion'', in which the crucified Christ towers
above the others, makes any other geometric emphasis of the compositional centre
superfluous. This is probably one of the reasons why the ``Crucifixion'' is always
painted against a geometrically inexpressive, flat Jerusalem wall.

\srcpage{209}

Since the main figures of Byzantine and Old Russian painting, unlike those of
Eastern miniatures, ``stand on the frame'', the depth of the depicted space can be
very small, in no way comparable to the depth of the action-filled space of the
miniatures discussed in this chapter. The effort in one case to bring the main
subject as close as possible to the viewer, and in the other to show some remove
from him---to give him the opportunity as it were to ``watch unobserved''---leads
to the result that in Byzantine and Old Russian art not only the depth of the
intended spatial layer is small, but (unlike Eastern miniatures) so is the distance
between the viewer and the first plane. Consequently, in accordance with the laws
of visual perception, various kinds of perspective transformation of the axonometric
basis of object depiction are here more distinctly expressed. Natural gentle reverse
perspective is more frequent, and the elongated rear legs of tables and seats play
an important role in the composition. In addition, transformations appear that are
not directly connected with the geometry of visual perception and that likewise
often emphasise reverse-perspective effects---entirely comprehensible in the context
of reverse perspective of natural origin. This seems a natural consequence of the
close artistic space of small depth, beginning behind the main figures.

Eastern miniatures, where the action takes place in sufficiently deep and distant
layers of space, are almost strictly axonometric.

Reverse perspective in Eastern miniatures, corresponding to natural visual
perception (described in Chapter~V), also occurs---less frequently and somewhat more
weakly than in Russian and Byzantine painting. As for effects of reverse perspective
(Chapter~VI), with the possible exception of those connected with the mobility of
the viewpoint (ill.\,94), they are virtually absent from Indian and Iranian
miniatures. There is here almost no hierarchical exaggeration of main figures; there
is no problem of ``non-occlusion'' since the depicted space is sufficiently large;
there are no compositional impulses connected with the presence in the image of a
large number of objects rendered in naturally-arising reverse perspective. As for
``informational rotations'' of planes, these give way to a preference for
technical-drawing methods (carpets and pools, for example, are shown exactly in
plan). Consequently, the artists of India and Iran followed visual perception of the
foreground more closely, and deviated from it less, than the icon painters. The
virtually sole departure from the laws of visual perception consists in the
transition to technical-drawing methods for rendering floors and other horizontal
planes---which in turn may require the introduction of geometrically contradictory
depictions. Both of these conventions certainly make miniatures difficult for the
modern viewer, who is unaccustomed to such devices.

\srcpage{210}

Both Indian and Iranian miniatures on the one hand, and the painting of Byzantium
and Old Russia on the other, are examples of the productive use of the system of
free perceptual perspective. These examples allow us to name several characteristic
features of such a perspective system.

The artist's freedom, as already stated, amounts to choosing the inevitable
distortions involved in rendering an individual object on the picture plane, while
striving in general to render each such object as far as possible in accordance with
visual perception. It is clear that these distortions are not chosen arbitrarily but
with regard to the construction of the image as a whole. In this connection the
artist is entitled to use multiple viewpoints, and where appropriate to incorporate
elements of technical drawing into the picture and to shift individual pictorial
elements in accordance with the logic of employing the method of geometrically
contradictory spatial rendering.

Availing himself of this freedom, the artist has the possibility of choosing the
geometric structure of his work in accordance with its functional characteristics.
In Chapter~VI it was already noted that such a functional approach can lead to
effects of reverse perspective. Let us consider this question in somewhat more
detail.

An Orthodox icon cannot be regarded merely as an illustration of a Biblical text or
of events described by church tradition---a kind of primer for the illiterate. Its
ritual function is to help the worshipper enter into communion with the heavenly
church, and this function of the icon could not but leave its mark on its artistic
characteristics. Byzantine and Old Russian religious painting possesses intensified
persuasive power, heightened appealingness. If one turns to the geometric problems
arising in this connection, one can assert in particular that such functionality is
capable of substantially affecting the construction of pictorial space in the icon.

As was noted above, the desire to draw the viewer's attention to the compositional
centre of the image led to the axes of architectural structures on the second plane
being directed toward that centre, giving rise to features of reverse perspective in
the overall space of the icon. But let us emphasise: the construction of icon space
on the scheme of reverse perspective not only draws the viewer's attention to the
compositional centre, but also alters the perception of that space.

If the linear perspective of Renaissance and modern art drew the viewer's gaze into
the depths of pictorial space, distancing the depicted from him; if the axonometry
of the medieval art of China or Japan made it possible to create paintings neutral
in this respect; then the reverse-perspective construction of space in icons produced
a sensation of the depicted space ``flowing toward'' the viewer---he as it were
became a participant in what was taking place in the icon.

\srcpage{211}

Let us examine this effect from the standpoint of the laws of visual perception.
The system of Renaissance linear perspective fully reproduces the monocular depth
cues discussed above, since it is based on the exact transmission of the retinal
image. This impels the brain of the viewer to unconsciously process the retinal
image in accordance with the size-constancy mechanism. Although such processing,
as was explained in Chapters~II and~III, cannot be complete, it nevertheless
generates a sensation that draws the viewer's gaze into the depth of the depicted
space.

If the space of the picture is constructed following the scheme of reverse
perspective, the overall effect is different. In the first place, such a scheme by
its very nature does not permit the rendering of deep spaces. Indeed, if an artist
were to attempt this, then, in progressively increasing the sizes of objects in the
image as depth increases, he would at once encounter the geometric fact that the
objects he is reproducing would, from a certain depth onward, simply not fit within
the bounds of the picture plane. But this near-obvious circumstance is not what is
important here. The point is that reverse perspective is conspicuously illogical if
one attempts to use it to construct the space of a picture---it is in glaring
contradiction with natural visual perception. This assertion requires extended
comment, since Chapter~V of this book was devoted to proving that gentle reverse
perspective is the norm of natural human visual perception. Both assertions are
simultaneously true, but they speak of different things.

It has been stressed many times above that natural perception in axonometry and
gentle reverse perspective is subject to two conditions: a small distance to the
depicted object and a small size of that object (more precisely, such a size that it
fits entirely within the zone of clear central vision). These conditions are usually
easily satisfied for the reproduction of individual objects. In other words,
axonometry and gentle reverse perspective are entirely natural for conveying the
volume and form of individual, not too large and sufficiently close objects
(footstools, seats, boxes, and the like). This aspect was in fact the primary one in
the chapters that linked reverse perspective to the regularities of visual
perception. If one speaks of the geometry of visual perception of large spaces, in
which the determining role is played by distant regions and regions covered by
indistinct peripheral vision, then here natural visual perception is characterised
by linear perspective very close to that developed in the Renaissance. Thus clearly
seen axonometry and gentle reverse perspective turn out to be, as it were, embedded
within a broad and deep space seen less clearly and according to the laws of linear
perspective.

\srcpage{212}

Medieval European art had no need to depict distant areas of space. The
compositions used were frequently characterised by a very shallow layer of space in
which figures were arranged as it were ``along the frame''. The small depth of the
rendered spatial layer allowed the principle of reverse perspective to be applied
not only to the depiction of individual objects but also to the rendering of the
whole of space, without encountering the difficulties mentioned somewhat above.

Here, of the five monocular depth cues named in Chapter~II, the two most
significant in the case under discussion should be recognised as the first two:
occlusion and the progressive diminution of the sizes of objects in the picture with
increasing depth. Generally speaking, the equally important third cue---aerial
perspective---is unable to manifest itself when depicting a space of small depth.

Comparing the first two depth cues, it is easy to notice their fundamental
difference. Occlusion (the masking of distant objects by closer ones) has to do with
the geometry of objective space; this cue divides the objects of the image by the
qualitative criterion of ``near--far'' and is fully capable of manifesting itself
not only in a drawing but even in a technical drawing (see Chapter~I). The second
cue (progressive diminution of the sizes of depicted objects with depth) has to do
only with perceptual space and is quantitative in character (from a photograph one
can establish the distances between objects).

In a picture built in the system of ordinary linear perspective, these two cues act
in concert and give the well-known sensation of depth familiar to everyone. In a
picture in which the entire space is constructed on the principles of reverse
perspective, an inevitable contradiction with the first depth cue---occlusion---
arises. As a result, the viewer's visual perception is placed in conditions where
two mutually exclusive pieces of information arrive through two channels
simultaneously, and it must somehow reconcile and interpret them.

\srcpage{213}

Experience in perceiving painting shows that no perspective devices are capable of
preventing the qualitatively correct perception of the mutual arrangement of
depicted objects, if the artist has used so powerful a depth cue as occlusion.
This is natural---occlusion is in some sense an absolute cue, directly connected
with the geometry of objective space. As for reverse perspective, it permits a
plausible subconscious interpretation. Gentle reverse perspective is natural when
rendering a small, close object, while architectural structures on the second plane
oriented toward the compositional centre can be perceived as buildings that are
genuinely oriented that way in objective space.

A depiction built on reverse perspective---in which not just one comparatively small
object is rendered according to this principle but the whole of the picture's space
is subordinated to it---evokes unconscious processes contrary to those that are
``natural'', analogous to those that give rise to various optical illusions. If
depth cues such as occlusion, along with equally weighty other depth cues, did not
as a whole provide a correct interpretation of the depicted, then a
reverse-perspective construction taken on its own would most likely be perceived
literally ``back to front''. Thus, a trapezoid widening toward the top, drawn on a
blank sheet of paper, would be perceived not as a floor (that is, in reverse
perspective) but rather as a ceiling---the large base of the trapezoid, understood
as a foreshortened rendering of a rectangle, would seem to be closer to the viewer
than its small base. But in reverse perspective the larger is the more distant, and
consequently the depiction of the far side would have ``moved toward'' the viewer.

As has just been said, other depth cues---and above all occlusion---do not permit
such ``liberties'' of perception; depicted objects do not ``flip''. But a certain
subconscious trace of these impulses, a certain resultant of all the depth-cue
information combined, leads to the result that instead of a sensation of depth
receding behind the picture plane, there arises a sensation of space peculiarly
``flowing toward'' the viewer without violation of the qualitative order of the
depicted objects in the ``near--far'' direction.

To strengthen this sensation, the artist not only builds the entire space on the
principle of reverse perspective (for example, by orienting the longitudinal axes of
architectural structures on the second plane accordingly), but also shows (contrary
to natural visual perception) the second-plane buildings in reverse perspective,
sharply intensifies the reverse perspective of first-plane objects (using a
conspicuously strong version instead of the natural gentle one), and so on.

The functional approach in medieval religious painting could stimulate recourse to
the system of reverse perspective not only in those cases where its appearance is
visually entirely natural, but also when it clearly contradicted the process of
visual spatial perception, in order to obtain the desired spatial illusions.

When the artist does not depart from the regularities of natural visual
perception---that is, when he applies gentle reverse perspective to individual
objects rather than to the whole of the space---none of the effects described above
are observed. This can be judged at least by the self-portrait of Kim Hong-do
(ill.\,39), in which gentle reverse perspective is applied only to the depiction of
a low table and the objects placed on it, while the space as a whole bears no
explicit trace of having been constructed according to the laws of reverse
perspective.

In the miniatures of India and Iran there is not even a hint of the perspective
exaggerations characteristic of many specimens of European medieval religious
painting.

In the light of what has been said, the dependence of spatial constructions on the
content of the artwork becomes obvious. This is what fundamentally distinguishes
free perceptual perspective from rigid systems (such as the linear one), in which
the geometry of spatial construction does not depend on what specific objects are
depicted, and where and for what purpose. In Eastern miniatures, although all the
elements of the methods of depiction are similar to those encountered in icon
painting, the spatial constructions as a whole are essentially different---owing to
the different emphasis placed on the common elements. This difference well
illustrates the freedom of perspective construction that has been spoken of many
times above.

Instead of unifying the picture as a whole by the introduction of a single
viewpoint (as in linear perspective), that unification is here achieved on a less
formal-geometric basis by various means of artistic expression.

Of course, the unification of a depiction into a coherent whole is achieved not by
the system of spatial constructions alone. Colour, rhythmic structures, symmetry,
ornamentality, and other analogous means play an important role here. All of this
lies outside the scope of the present book, which is devoted to the comparatively
narrow question of the geometry of spatial constructions.
